\documentclass{article}
\usepackage[utf8]{inputenc}
\usepackage[spanish]{babel}
\usepackage{listings}
\usepackage{xcolor}
\usepackage{geometry}
\usepackage{amsmath}
\geometry{a4paper, margin=1in}
\definecolor{codegreen}{rgb}{0,0.6,0}
\definecolor{codegray}{rgb}{0.5,0.5,0.5}
\definecolor{codepurple}{rgb}{0.58,0,0.82}
\definecolor{backcolour}{rgb}{0.95,0.95,0.92}
\lstset{backgroundcolor=\color{backcolour},commentstyle=\color{codegreen},keywordstyle=\color{magenta},numberstyle=\tiny\color{codegray},stringstyle=\color{codepurple},basicstyle=\ttfamily\small,breakatwhitespace=false,breaklines=true,captionpos=b,keepspaces=true,numbers=left,numbersep=5pt,showspaces=false,showstringspaces=false,showtabs=false,tabsize=2,language=C++}
\newcommand{\quees}[1]{\section*{🤔 ¿QUÉ ES?}\hrulefill \\ \#1}
\newcommand{\dichodeotraforma}[1]{\section*{💬 DICHO DE OTRA FORMA}\hrulefill \\ \#1}
\newcommand{\diagrama}[1]{\section*{🎨 DIAGRAMA}\hrulefill \\ \#1}
\newcommand{\comofunciona}[1]{\section*{⚙️ ¿CÓMO FUNCIONA?}\hrulefill \\ \#1}
\newcommand{\complejidad}[1]{\section*{⏱️ COMPLEJIDAD (Big-O)}\hrulefill \\ \#1}
\newcommand{\entrevistas}[1]{\section*{🎯 EN ENTREVISTAS}\hrulefill \\ \#1}
\newcommand{\resumen}[1]{\section*{📋 RESUMEN RÁPIDO}\hrulefill \\ \#1}
\begin{document}
\title{Two Pointers (Dos Punteros)}
\date{}
\maketitle
\quees{Dos personas en extremos opuestos de un pasillo 🚶‍♂️...🚶‍♀️ caminando 
hacia el centro para encontrarse. O un corredor lento y uno rápido 
en la misma pista 🏃‍♂️💨...🏃‍♂️.

En lugar de usar un solo índice para recorrer un array (un loop 
simple), usas DOS. Estos punteros se mueven de forma coordinada
para resolver problemas de búsqueda o comparación mucho más rápido
que la fuerza bruta.

Es uno de los patrones más útiles para entrevistas técnicas porque 
suele reducir soluciones de O(n^2) a O(n).}
\dichodeotraforma{El patrón Two Pointers utiliza dos variables de índice para recorrer
una estructura de datos lineal (array, string, linked list). 

Existen dos variantes principales:
1. Opposite Ends: Los punteros empiezan en los extremos y se mueven
   hacia el centro. (Ej: Two Sum en array ordenado, Palíndromos).
2. Fast \& Slow: Un puntero avanza más rápido que el otro. (Ej: Detectar
   ciclos en Linked List, encontrar el medio).}
\diagrama{\begin{verbatim}
Problema: Encontrar par que sume 10 en array ordenado.
[ 1, 3, 4, 6, 8, 9 ]
  L              R   -> 1 + 9 = 10. ¡ENCONTRADO! 🎉

Si la suma fuera menor a 10, moveríamos L a la derecha.
Si la suma fuera mayor a 10, moveríamos R a la izquierda.
\end{verbatim}}
\comofunciona{Opposite Ends:
1. Inicializas `left = 0` y `right = n - 1`.
2. Mientras `left < right`:
   - Evalúas la condición con `arr[left]` y `arr[right]`.
   - Si se cumple la meta, terminas.
   - Si no, decides qué puntero mover según la lógica del problema.

Fast \& Slow (Floyd's Algorithm):
1. Inicializas `slow = start` y `fast = start`.
2. En cada paso: `slow` avanza 1, `fast` avanza 2.
3. Si `fast` alcanza el final, no hay ciclo.
4. Si `slow == fast`, ¡hay un ciclo!}
\section*{📝 PSEUDOCÓDIGO}
\hrulefill
\begin{verbatim}
// Two Sum en Array Ordenado
función tieneSuma(arr, meta):
    izq = 0
    der = arr.tamaño - 1
    MIENTRAS izq < der:
        sumaActual = arr[izq] + arr[der]
        SI sumaActual == meta: return VERDADERO
        SI sumaActual < meta:  izq++
        SINO:                  der--
    return FALSO
\end{verbatim}
\section*{💻 CÓDIGO C++}
\hrulefill
\begin{lstlisting}[language=C++]
#include <iostream>
#include <vector>
#include <string>
#include <algorithm>
using namespace std;

// Verificar si un string es palíndromo -> O(n) tiempo, O(1) espacio
bool isPalindrome(string s) {
    int left = 0;
    int right = s.length() - 1;

    while (left < right) {
        if (s[left] != s[right]) return false;
        left++;
        right--;
    }
    return true;
}

int main() {
    string word = "reconocer";
    cout << "¿Es palíndromo? " << (isPalindrome(word) ? "SÍ" : "NO") << endl;
    
    return 0;
}
\end{lstlisting}
\complejidad{Caso        │ Tiempo        │ Espacio
  ────────────┼───────────────┼────────────────────────────────
  Two Pointers│ O(n)          │ O(1)
  Fuerza Bruta│ O(n^2)         │ O(1)

* La magia del patrón es que evitas el segundo loop anidado.}
\entrevistas{}
\resumen{• Dos índices moviéndose coordinados.
• Reduce O(n^2) a O(n).
• Memoria constante O(1) casi siempre.
• Variantes: Opuestos (extremos → centro) y Fast/Slow.
• Requisito común: Array ordenado.
• Indispensable para optimizar búsquedas de pares o ciclos.
════════════════════════════════════════════════════════════════}
\end{document}